\section{Determinants: Factor Hunting, Laplace Expansion, and Desnanot--Jacobi}

\subsection{Factor hunting and the Vandermonde determinant}

\begin{theorem}
[Vandermonde determinant]
\label{thm.det.vander}
\lean{AlgebraicCombinatorics.Determinants.vandermonde_det_a, AlgebraicCombinatorics.Determinants.vandermonde_det_b, AlgebraicCombinatorics.Determinants.vandermonde_det_c, AlgebraicCombinatorics.Determinants.vandermonde_det_d}
\leantarget
\leanok
Let $n \in \mathbb{N}$. Let $a_1, a_2, \ldots, a_n$ be $n$ elements of $K$. Then:

\textbf{(a)} We have
\[
  \det\left( \left( a_i^{n-j} \right)_{1 \leq i \leq n,\; 1 \leq j \leq n} \right)
  = \prod_{1 \leq i < j \leq n} (a_i - a_j).
\]

\textbf{(b)} We have
\[
  \det\left( \left( a_j^{n-i} \right)_{1 \leq i \leq n,\; 1 \leq j \leq n} \right)
  = \prod_{1 \leq i < j \leq n} (a_i - a_j).
\]

\textbf{(c)} We have
\[
  \det\left( \left( a_i^{j-1} \right)_{1 \leq i \leq n,\; 1 \leq j \leq n} \right)
  = \prod_{1 \leq j < i \leq n} (a_i - a_j).
\]

\textbf{(d)} We have
\[
  \det\left( \left( a_j^{i-1} \right)_{1 \leq i \leq n,\; 1 \leq j \leq n} \right)
  = \prod_{1 \leq j < i \leq n} (a_i - a_j).
\]
\end{theorem}

\begin{proof}
\textbf{(a)} Follows from Lemma~\ref{lem.det.vander.a.pol} by substituting
$a_1, a_2, \ldots, a_n$ for the indeterminates $x_1, x_2, \ldots, x_n$.
More precisely, the equality
\[
  \det\left( \left( a_i^{n-j} \right)_{1 \leq i \leq n,\; 1 \leq j \leq n} \right)
  = \prod_{1 \leq i < j \leq n} (a_i - a_j)
\]
follows from the polynomial identity~\eqref{eq.lem.det.vander.a.pol.1} by
applying the evaluation homomorphism $x_i \mapsto a_i$.

\textbf{(b)} The matrix $\left( a_j^{n-i} \right)_{1 \leq i \leq n,\; 1 \leq j \leq n}$
is the transpose of the matrix $\left( a_i^{n-j} \right)_{1 \leq i \leq n,\; 1 \leq j \leq n}$.
Thus part~(b) follows from part~(a) using $\det(A^T) = \det A$.

\textbf{(c)} The matrix $\left( a_i^{j-1} \right)$ is obtained from
$\left( a_i^{n-j} \right)$ by reversing the column order.
Let $\tau \in S_n$ send each $j$ to $n+1-j$. Then
\[
  \det\left( \left( a_i^{j-1} \right) \right)
  = (-1)^\tau \cdot \det\left( \left( a_i^{n-j} \right) \right)
  = (-1)^\tau \cdot \prod_{1 \leq i < j \leq n} (a_i - a_j).
\]
Since each factor $(a_i - a_j)$ gets negated by $(-1)^\tau$, the result
becomes $\prod_{1 \leq j < i \leq n} (a_i - a_j)$.

\textbf{(d)} Follows from part~(c) using $\det(A^T) = \det A$.
\end{proof}

\begin{lemma}
\label{lem.det.vander.a.pol}
\lean{AlgebraicCombinatorics.Determinants.vandermonde_det_poly'}
\leantarget
\leanok
Let $n \in \mathbb{N}$. Consider the polynomial ring
$\mathbb{Z}[x_1, x_2, \ldots, x_n]$ in $n$ indeterminates
$x_1, x_2, \ldots, x_n$ with integer coefficients. In this ring, we have
\begin{equation}
  \det\left( \left( x_i^{n-j} \right)_{1 \leq i \leq n,\; 1 \leq j \leq n} \right)
  = \prod_{1 \leq i < j \leq n} (x_i - x_j).
  \label{eq.lem.det.vander.a.pol.1}
\end{equation}
\end{lemma}

\begin{proof}
Set
\[
  f := \det\left( \left( x_i^{n-j} \right)_{1 \leq i \leq n,\; 1 \leq j \leq n} \right)
  \quad \text{and} \quad
  g := \prod_{1 \leq i < j \leq n} (x_i - x_j).
\]
We must prove that $f = g$.

By the definition of the determinant,
\[
  f = \sum_{\sigma \in S_n} (-1)^\sigma
      x_1^{n - \sigma(1)} x_2^{n - \sigma(2)} \cdots x_n^{n - \sigma(n)},
\]
which is a homogeneous polynomial of degree $\frac{n(n-1)}{2}$.
The monomial $x_1^{n-1} x_2^{n-2} \cdots x_n^0$ appears with coefficient~$1$
(from the identity permutation $\sigma = \mathrm{id}$).

For any $u < v$ in $[n]$, setting $x_u = x_v$ makes two rows of the matrix
identical, so $f$ vanishes. By the root factoring-off theorem (viewing $f$ as a
polynomial in $x_u$ over
$\mathbb{Z}[x_1, \ldots, \widehat{x_u}, \ldots, x_n]$), we conclude that
$x_u - x_v$ divides $f$. Since the ring
$\mathbb{Z}[x_1, \ldots, x_n]$ is a UFD and the linear factors $x_u - x_v$
are irreducible and mutually non-associate, the product
$g = \prod_{1 \leq i < j \leq n} (x_i - x_j)$ divides $f$. Thus $f = gq$
for some $q \in \mathbb{Z}[x_1, \ldots, x_n]$.

Since $\deg f \leq \frac{n(n-1)}{2} = \deg g$, we get $\deg q \leq 0$, so
$q \in \mathbb{Z}$. Comparing the coefficient of $x_1^{n-1} x_2^{n-2} \cdots x_n^0$
on both sides gives $1 = q \cdot 1$, hence $q = 1$ and $f = g$.

The formalization relates our matrix convention to Mathlib's via column reversal,
and uses the fact that swapping the order of subtraction in each factor introduces
a sign that cancels with the sign from the column permutation.
\end{proof}

\begin{proposition}
\label{prop.det.(xi+yj)n-1}
\lean{AlgebraicCombinatorics.Determinants.det_sum_pow}
\leantarget
\leanok
Let $n \in \mathbb{N}$. Let $x_1, x_2, \ldots, x_n \in K$ and
$y_1, y_2, \ldots, y_n \in K$. Then,
\begin{align*}
  & \det\left( \left( (x_i + y_j)^{n-1} \right)_{1 \leq i \leq n,\; 1 \leq j \leq n} \right) \\
  &= \left( \prod_{k=0}^{n-1} \binom{n-1}{k} \right)
     \left( \prod_{1 \leq i < j \leq n} (x_i - x_j) \right)
     \left( \prod_{1 \leq i < j \leq n} (y_j - y_i) \right).
\end{align*}
\end{proposition}

\begin{proof}
Let $C = \left( (x_i + y_j)^{n-1} \right)_{1 \leq i \leq n,\; 1 \leq j \leq n}$.
By the binomial formula, for any $i, j \in [n]$,
\[
  C_{i,j} = (x_i + y_j)^{n-1}
           = \sum_{k=1}^{n} \binom{n-1}{k-1} x_i^{k-1} y_j^{n-k}.
\]
Define two $n \times n$-matrices
\[
  P := \left( \binom{n-1}{k-1} x_i^{k-1} \right)_{1 \leq i \leq n,\; 1 \leq k \leq n}
  \quad \text{and} \quad
  Q := \left( y_j^{n-k} \right)_{1 \leq k \leq n,\; 1 \leq j \leq n}.
\]
Then $C = PQ$, so $\det C = \det P \cdot \det Q$.

From $Q = \left( y_j^{n-i} \right)$, Theorem~\ref{thm.det.vander}~(b) gives
$\det Q = \prod_{1 \leq i < j \leq n} (y_i - y_j)$.

From $P = \left( \binom{n-1}{j-1} x_i^{j-1} \right)$, factoring out the
binomial coefficients from each column and applying
Theorem~\ref{thm.det.vander}~(c) gives
\[
  \det P = \left( \prod_{k=0}^{n-1} \binom{n-1}{k} \right)
           \cdot \prod_{1 \leq j < i \leq n} (x_i - x_j)
         = \left( \prod_{k=0}^{n-1} \binom{n-1}{k} \right)
           \cdot \prod_{1 \leq i < j \leq n} (x_j - x_i).
\]

Hence,
\begin{align*}
  \det C &= \det P \cdot \det Q \\
         &= \left( \prod_{k=0}^{n-1} \binom{n-1}{k} \right)
            \cdot \prod_{1 \leq i < j \leq n}
            \underbrace{(x_j - x_i)(y_i - y_j)}_{= (x_i - x_j)(y_j - y_i)} \\
         &= \left( \prod_{k=0}^{n-1} \binom{n-1}{k} \right)
            \left( \prod_{1 \leq i < j \leq n} (x_i - x_j) \right)
            \left( \prod_{1 \leq i < j \leq n} (y_j - y_i) \right).
\end{align*}
\end{proof}

\subsection{Laplace expansion}

\begin{convention}
\label{conv.mat.tilde}
\lean{AlgebraicCombinatorics.Determinants.submatrixRemove}
\leanok
Let $n \in \mathbb{N}$. Let $A$ be an $n \times n$-matrix.
Let $i, j \in [n]$. Then, we set
\[
  A_{\sim i, \sim j} := \operatorname{sub}_{[n] \setminus \{i\}}^{[n] \setminus \{j\}} A.
\]
This is the $(n-1) \times (n-1)$-matrix obtained from $A$ by removing its
$i$-th row and its $j$-th column.
\end{convention}

\begin{theorem}
[Laplace expansion]
\label{thm.det.laplace}
\lean{AlgebraicCombinatorics.Determinants.det_laplace_row, AlgebraicCombinatorics.Determinants.det_laplace_col}
\leantarget
\leanok
Let $n \in \mathbb{N}$. Let $A \in K^{n \times n}$ be an $n \times n$-matrix.

\textbf{(a)} For every $p \in [n]$, we have
\[
  \det A = \sum_{q=1}^{n} (-1)^{p+q} A_{p,q} \det(A_{\sim p, \sim q}).
\]

\textbf{(b)} For every $q \in [n]$, we have
\[
  \det A = \sum_{p=1}^{n} (-1)^{p+q} A_{p,q} \det(A_{\sim p, \sim q}).
\]
\end{theorem}

\begin{proof}
See \cite[Theorem 6.82]{detnotes} or
\cite[Lemma 4.7.7]{Ford21} or \cite[5.8 and 5.8']{Laue-det} or
\cite[Proposition B.24 and Proposition B.25]{Strick13} or \cite[Theorem
9.48]{Loehr-BC}.

Part~(a) is Mathlib's \texttt{Matrix.det\_succ\_row}; part~(b) is
\texttt{Matrix.det\_succ\_column}.
\end{proof}

\begin{proposition}
\label{prop.det.laplace.0}
\lean{AlgebraicCombinatorics.Determinants.det_laplace_row_zero, AlgebraicCombinatorics.Determinants.det_laplace_col_zero}
\leantarget
\leanok
Let $n \in \mathbb{N}$. Let $A \in K^{n \times n}$ be an $n \times n$-matrix.
Let $r \in [n]$.

\textbf{(a)} For every $p \in [n]$ satisfying $p \neq r$, we have
\[
  0 = \sum_{q=1}^{n} (-1)^{p+q} A_{r,q} \det(A_{\sim p, \sim q}).
\]

\textbf{(b)} For every $q \in [n]$ satisfying $q \neq r$, we have
\[
  0 = \sum_{p=1}^{n} (-1)^{p+q} A_{p,r} \det(A_{\sim p, \sim q}).
\]
\end{proposition}

\begin{proof}
\textbf{(a)} Let $p \in [n]$ satisfy $p \neq r$. Let $C$ be the matrix $A$
with its $p$-th row replaced by its $r$-th row. Then $C$ has two equal rows,
so $\det C = 0$. On the other hand, expanding $\det C$ along the $p$-th row
(using Theorem~\ref{thm.det.laplace}~(a) applied to $C$) yields
\[
  \det C = \sum_{q=1}^{n} (-1)^{p+q} \underbrace{C_{p,q}}_{= A_{r,q}}
           \det\underbrace{(C_{\sim p, \sim q})}_{= A_{\sim p, \sim q}}
         = \sum_{q=1}^{n} (-1)^{p+q} A_{r,q} \det(A_{\sim p, \sim q}).
\]
Comparing gives the claim. Part~(b) follows similarly.

The formalization uses the adjugate: $(-1)^{p+q} \det(A_{\sim p, \sim q})
= (\operatorname{adj} A)_{q,p}$, so the sum becomes
$\sum_q A_{r,q} (\operatorname{adj} A)_{q,p} = (A \cdot \operatorname{adj} A)_{r,p}
= (\det A) \cdot I_{r,p} = 0$ when $r \neq p$.
\end{proof}

\begin{definition}
[Adjugate matrix]
\label{def.det.adj}
\lean{AlgebraicCombinatorics.Determinants.adjugateMat, AlgebraicCombinatorics.Determinants.adjugateMat_eq_adjugate}
\leantarget
\leanok
Let $n \in \mathbb{N}$. Let $A \in K^{n \times n}$ be an $n \times n$-matrix.
We define a new $n \times n$-matrix $\operatorname{adj} A \in K^{n \times n}$ by
\[
  \operatorname{adj} A
  = \left( (-1)^{i+j} \det(A_{\sim j, \sim i}) \right)_{1 \leq i \leq n,\; 1 \leq j \leq n}.
\]
This matrix $\operatorname{adj} A$ is called the \emph{adjugate} (or
\emph{classical adjoint}) of the matrix $A$.
Note the index swap: the $(i,j)$ entry involves removing row $j$ and column $i$.

The formalization shows that this definition equals Mathlib's
\texttt{Matrix.adjugate}, which is defined via
$(\operatorname{adj} A)_{i,j} = \det(A')$
where $A'$ is the matrix $A$ with row $j$ replaced by the $i$-th standard basis vector.
\end{definition}

\begin{theorem}
\label{thm.det.adj.inverse}
\lean{AlgebraicCombinatorics.Determinants.mul_adjugate', AlgebraicCombinatorics.Determinants.adjugate_mul'}
\leantarget
\leanok
Let $n \in \mathbb{N}$. Let $A \in K^{n \times n}$ be an $n \times n$-matrix.
Let $I_n$ denote the $n \times n$ identity matrix. Then,
\[
  A \cdot (\operatorname{adj} A) = (\operatorname{adj} A) \cdot A
  = (\det A) \cdot I_n.
\]
\end{theorem}

\begin{proof}
See \cite[Theorem 6.100]{detnotes} or \cite[Lemma 4.7.9]{Ford21} or
\cite[Theorem 9.50]{Loehr-BC} or \cite[Proposition B.28]{Strick13}.

To show $A \cdot (\operatorname{adj} A) = (\det A) \cdot I_n$, it suffices to
check that the $(i,j)$-th entry of $A \cdot (\operatorname{adj} A)$ equals
$\det A$ when $i = j$ and equals $0$ otherwise.

Case $i = j$: The $(i,i)$ entry is
$\sum_k A_{i,k} (\operatorname{adj} A)_{k,i}
= \sum_k (-1)^{i+k} A_{i,k} \det(A_{\sim i, \sim k}) = \det A$
by Laplace expansion along row $i$ (Theorem~\ref{thm.det.laplace}~(a)).

Case $i \neq j$: The $(i,j)$ entry is
$\sum_k A_{i,k} (\operatorname{adj} A)_{k,j}
= \sum_k (-1)^{j+k} A_{i,k} \det(A_{\sim j, \sim k}) = 0$
by Proposition~\ref{prop.det.laplace.0}~(a) (using row~$i$ entries with row~$j$ cofactors).

Similarly, $(\operatorname{adj} A) \cdot A = (\det A) \cdot I_n$ follows
using column versions.
\end{proof}

\subsection{Laplace expansion along multiple rows}

\begin{convention}
\label{conv.det.sum-compl}
For any subset $I$ of $[n]$, we write $\widetilde{I} = [n] \setminus I$ for
its complement, and $\operatorname{sum} S = \sum_{s \in S} s$ for any finite
set $S$ of integers.
\end{convention}

\begin{theorem}
[Laplace expansion along multiple rows/columns]
\label{thm.det.laplace-multi}
\lean{AlgebraicCombinatorics.Determinants.det_laplace_multi_row, AlgebraicCombinatorics.Determinants.det_laplace_multi_col}
\leantarget
\leanok
Let $n \in \mathbb{N}$. Let $A \in K^{n \times n}$ be an $n \times n$-matrix.

\textbf{(a)} For every subset $P$ of $[n]$, we have
\[
  \det A = \sum_{\substack{Q \subseteq [n] \\ |Q| = |P|}}
           (-1)^{\operatorname{sum} P + \operatorname{sum} Q}
           \det\left( \operatorname{sub}_P^Q A \right)
           \det\left( \operatorname{sub}_{\widetilde{P}}^{\widetilde{Q}} A \right).
\]

\textbf{(b)} For every subset $Q$ of $[n]$, we have
\[
  \det A = \sum_{\substack{P \subseteq [n] \\ |P| = |Q|}}
           (-1)^{\operatorname{sum} P + \operatorname{sum} Q}
           \det\left( \operatorname{sub}_P^Q A \right)
           \det\left( \operatorname{sub}_{\widetilde{P}}^{\widetilde{Q}} A \right).
\]
\end{theorem}

\begin{proof}
See \cite[Theorem 6.156]{detnotes}.

The proof partitions the symmetric group based on how permutations map $P$
to subsets $Q$ of the same size. For each such $Q$, the permutations
$\sigma$ with $\sigma(P) = Q$ decompose as a bijection $\tau : P \to Q$
(contributing to $\det(\operatorname{sub}_P^Q A)$) and a bijection
$\rho : \widetilde{P} \to \widetilde{Q}$
(contributing to $\det(\operatorname{sub}_{\widetilde{P}}^{\widetilde{Q}} A)$).
The sign decomposes as
$(-1)^\sigma = (-1)^{\operatorname{sum} P + \operatorname{sum} Q}
(-1)^\tau (-1)^\rho$.

Part~(b) follows from part~(a) applied to $A^T$ using $\det(A^T) = \det A$.
\end{proof}

\subsection{Desnanot--Jacobi and Dodgson condensation}

\begin{theorem}
[Desnanot--Jacobi formula, take 1]
\label{thm.det.des-jac-1}
\lean{AlgebraicCombinatorics.Determinants.desnanot_jacobi, AlgebraicCombinatorics.Determinants.desnanot_jacobi_det2}
\leantarget
\leanok
Let $n \in \mathbb{N}$ be such that $n \geq 2$. Let $A \in K^{n \times n}$ be
an $n \times n$-matrix.

Let $A'$ be the $(n-2) \times (n-2)$-matrix
\[
  \operatorname{sub}_{\{2, 3, \ldots, n-1\}}^{\{2, 3, \ldots, n-1\}} A
  = \left( A_{i+1, j+1} \right)_{1 \leq i \leq n-2,\; 1 \leq j \leq n-2}.
\]
(This is what remains of $A$ when we remove the first row, the last row, the
first column and the last column.) Then,
\begin{align*}
  \det A \cdot \det(A')
  &= \det(A_{\sim 1, \sim 1}) \cdot \det(A_{\sim n, \sim n})
     - \det(A_{\sim 1, \sim n}) \cdot \det(A_{\sim n, \sim 1}) \\
  &= \det \begin{pmatrix}
       \det(A_{\sim 1, \sim 1}) & \det(A_{\sim 1, \sim n}) \\
       \det(A_{\sim n, \sim 1}) & \det(A_{\sim n, \sim n})
     \end{pmatrix}.
\end{align*}
\end{theorem}

\begin{proof}
See \cite[Proposition 6.122]{detnotes} or \cite[\S 3.5, proof of Theorem 3.12]{Bresso99} or
\cite{zeilberger-twotime}.

The proof uses the adjugate matrix.

The right-hand side equals
$(\operatorname{adj} A)_{0,0} \cdot (\operatorname{adj} A)_{n-1,n-1}
- (\operatorname{adj} A)_{0,n-1} \cdot (\operatorname{adj} A)_{n-1,0}$,
which is the $2 \times 2$ determinant of the corner submatrix of
$\operatorname{adj} A$.

By Jacobi's complementary minor theorem for the $2 \times 2$ case (a special
case proved independently), this $2 \times 2$ determinant of the adjugate
equals $\det A \cdot \det(A')$, giving the left-hand side.

The proof of the special case uses the
polynomial ring approach: reduce to the generic matrix over
the multivariate polynomial ring, then work in its field of fractions where the matrix
is invertible.
\end{proof}

\begin{theorem}
[Cauchy determinant]
\label{thm.det.cauchy}
\lean{AlgebraicCombinatorics.Determinants.cauchy_det}
\leantarget
\leanok
Let $n \in \mathbb{N}$. Let $x_1, x_2, \ldots, x_n$ be $n$ elements of $K$.
Let $y_1, y_2, \ldots, y_n$ be $n$ elements of $K$.
Assume that $x_i + y_j$ is invertible in $K$ for each $(i, j) \in [n]^2$. Then,
\[
  \det\left( \left( \frac{1}{x_i + y_j} \right)_{1 \leq i \leq n,\; 1 \leq j \leq n}
  \right)
  = \frac{\prod_{1 \leq i < j \leq n} \left( (x_i - x_j)(y_i - y_j) \right)}{
          \prod_{(i,j) \in [n]^2} (x_i + y_j)}.
\]
\end{theorem}

\begin{proof}
The proof uses the polynomial version of the identity (factor hunting).

\textbf{Step 1: Clear denominators.}
Multiply both sides by $\prod_{i,j}(x_i + y_j)$ to obtain the ``cleared''
polynomial identity:
\[
  \det\!\left(\left(\prod_{k \neq j}(x_i + y_k)\right)_{i,j}\right)
  = \prod_{1 \leq i < j \leq n}\left((x_i - x_j)(y_i - y_j)\right).
\]

\textbf{Step 2: Base cases.}
For $n = 0, 1, 2, 3$: verified by direct computation.

\textbf{Step 3: Factor hunting} ($n \geq 4$):
The determinant on the left is a polynomial in $x_1, \ldots, x_n, y_1, \ldots, y_n$.
Setting $x_i = x_j$ (for $i \neq j$) makes two rows identical, so
$(x_i - x_j)$ divides the determinant.
Similarly, $(y_i - y_j)$ divides the determinant for $i \neq j$.
Since these linear factors are pairwise coprime irreducibles in the UFD
$\mathbb{Z}[x_1, \ldots, x_n, y_1, \ldots, y_n]$, their product divides
the determinant.

\textbf{Step 4: Degree comparison.}
Both the determinant and the product have total degree $n(n-1)$, so the
quotient is a constant. Evaluating at a specific point shows the constant
is~$1$.

\textbf{Step 5: Recover the Cauchy formula.}
Dividing the cleared identity by $\prod_{i,j}(x_i + y_j)$ yields the
Cauchy determinant formula.
\end{proof}

\begin{lemma}
\label{lem.det.cauchy-poly}
\lean{AlgebraicCombinatorics.Determinants.cauchy_det_poly}
\leanhelper
\leanok
Let $n \in \mathbb{N}$. Let $x_1, \ldots, x_n \in R$ and $y_1, \ldots, y_n \in R$
be elements of a commutative ring $R$. Then
\[
  \det\left( \left( \prod_{k \neq j} (x_i + y_k) \right)_{1 \leq i \leq n,\; 1 \leq j \leq n} \right)
  = \prod_{1 \leq i < j \leq n} (x_i - x_j)(y_i - y_j).
\]
This is equivalent to the Cauchy determinant formula after dividing both sides
by $\prod_{(i,j) \in [n]^2} (x_i + y_j)$.
\end{lemma}

\begin{proof}
The proof uses factor hunting in the polynomial ring. Both sides are polynomials
of degree $n(n-1)$ in the $x_i$'s and $y_j$'s. The left-hand side vanishes
when $x_u = x_v$ (two rows become identical) and when $y_u = y_v$ (two columns
become proportional). Thus the product
$\prod_{i < j} (x_i - x_j)(y_i - y_j)$ divides the left-hand side.
Degree comparison shows the quotient is a constant, and coefficient
comparison shows the constant is~$1$.
\end{proof}

\subsection{Generalized Desnanot--Jacobi}

\begin{theorem}
\label{thm.det.des-jac-2}
\lean{AlgebraicCombinatorics.Determinants.desnanot_jacobi_general}
\leantarget
\leanok
Let $n \in \mathbb{N}$ be such that $n \geq 2$. Let $p, q, u, v$ be four
elements of $[n]$ such that $p < q$ and $u < v$. Let $A$ be an
$n \times n$-matrix. Then,
\begin{align*}
  & \det A \cdot \det\left( \operatorname{sub}_{[n] \setminus \{p, q\}}^{[n] \setminus \{u, v\}} A \right) \\
  &= \det(A_{\sim p, \sim u}) \cdot \det(A_{\sim q, \sim v})
     - \det(A_{\sim p, \sim v}) \cdot \det(A_{\sim q, \sim u}).
\end{align*}
\end{theorem}

\begin{proof}
See \cite[Theorem 6.126]{detnotes}.

The proof combines two identities:
\begin{enumerate}
\item \textbf{Adjugate identity} (Lemma~\ref{lem.det.adj-corners}, generalized): The $2 \times 2$
      determinant
      $(\operatorname{adj} A)_{u,p} (\operatorname{adj} A)_{v,q}
       - (\operatorname{adj} A)_{u,q} (\operatorname{adj} A)_{v,p}$
      equals $(-1)^{p+u+q+v}$ times the product-minus-product expression.

\item \textbf{Jacobi $2 \times 2$ identity} (Theorem~\ref{thm.det.jacobi-complement} for $|P|=|Q|=2$): The same
      $2 \times 2$ determinant of the adjugate equals
      $(-1)^{p+u+q+v} \det A \cdot \det(\operatorname{sub}_{[n] \setminus \{p,q\}}^{[n] \setminus \{u,v\}} A)$.
\end{enumerate}
Equating the two expressions and canceling $(-1)^{p+u+q+v}$ (which squares to~$1$)
yields the result.

Theorem~\ref{thm.det.des-jac-1} is the particular case $p = 1$, $q = n$,
$u = 1$, $v = n$.
\end{proof}

\subsection{Jacobi's complementary minor theorem}

\begin{theorem}
[Jacobi's complementary minor theorem for adjugates]
\label{thm.det.jacobi-complement}
\lean{AlgebraicCombinatorics.Determinants.jacobi_complementary_minor}
\leantarget
\leanok
Let $n \in \mathbb{N}$. For any subset $I$ of $[n]$, let $\widetilde{I}$
denote the complement $[n] \setminus I$. Set
$\operatorname{sum} S = \sum_{s \in S} s$ for any finite set $S$ of integers.

Let $A \in K^{n \times n}$ be an $n \times n$-matrix. Let $P$ and $Q$ be two
subsets of $[n]$ such that $|P| = |Q| \geq 1$. Then,
\[
  \det\left( \operatorname{sub}_P^Q (\operatorname{adj} A) \right)
  = (-1)^{\operatorname{sum} P + \operatorname{sum} Q}
    (\det A)^{|Q| - 1}
    \det\left( \operatorname{sub}_{\widetilde{Q}}^{\widetilde{P}} A \right).
\]
\end{theorem}

\begin{proof}
The proof uses the polynomial ring approach (following Prasolov's
\emph{Problems and Theorems in Linear Algebra}).

\textbf{Step 1: Reduce to generic matrix.}
Let $A'$ be the generic $n \times n$ matrix with polynomial entries in
$\mathbb{Z}[(x_{i,j})_{1 \leq i,j \leq n}]$.
The identity for arbitrary $A$ over any commutative ring $R$ follows by
applying the evaluation homomorphism $x_{i,j} \mapsto A_{i,j}$.

\textbf{Step 2: Work in the field of fractions.}
Since $\det(A')$ is a nonzero polynomial, we can embed $A'$ into the field of
fractions $K$ of the polynomial ring, where $A'$ becomes
invertible.

\textbf{Step 3: Use $\operatorname{adj}(A) = \det(A) \cdot A^{-1}$.}
For invertible $A$ with $|P| = |Q| = k$:
\[
  \det(\operatorname{sub}_P^Q(\operatorname{adj} A))
  = (\det A)^k \cdot \det(\operatorname{sub}_P^Q(A^{-1})).
\]

\textbf{Step 4: Complementary minor theorem for inverse.}
The classical result for invertible $A$ states:
\[
  \det(\operatorname{sub}_P^Q(A^{-1}))
  = (-1)^{\operatorname{sum} P + \operatorname{sum} Q}
    \frac{\det(\operatorname{sub}_{\widetilde{Q}}^{\widetilde{P}} A)}{\det A}.
\]
This is proved using block matrix decomposition.

\textbf{Step 5: Combine.}
$\det(\operatorname{sub}_P^Q(\operatorname{adj} A))
= (\det A)^k \cdot (-1)^{\operatorname{sum} P + \operatorname{sum} Q}
  \cdot \det(\operatorname{sub}_{\widetilde{Q}}^{\widetilde{P}} A) / \det A
= (-1)^{\operatorname{sum} P + \operatorname{sum} Q}
  (\det A)^{k-1} \det(\operatorname{sub}_{\widetilde{Q}}^{\widetilde{P}} A)$.

\textbf{Step 6: Pull back.}
Since the identity holds in $K$ and the canonical map from the polynomial ring
to $K$ is injective, the identity holds in the polynomial ring, and hence for
all matrices $A$ over any commutative ring.

Theorem~\ref{thm.det.des-jac-2} is the particular case of
Theorem~\ref{thm.det.jacobi-complement} for $P = \{u, v\}$ and $Q = \{p, q\}$.
Theorem~\ref{thm.det.des-jac-1} is, of course, the particular case of
Theorem~\ref{thm.det.des-jac-2} for $p = 1$, $q = n$, $u = 1$, $v = n$.
\end{proof}

\subsection{Helpers for the Vandermonde determinant proof}

\begin{lemma}
\label{lem.det.vander.prod-Ioi-neg-sub}
\lean{AlgebraicCombinatorics.Determinants.prod_Ioi_neg_sub'}
\leanhelper
\leanok
Let $v_1, \ldots, v_m$ be elements of a commutative ring $S$. Then
\[
  \prod_{1 \leq i < j \leq m} (v_i - v_j)
  = (-1)^{m(m-1)/2} \prod_{1 \leq i < j \leq m} (v_j - v_i).
\]
\end{lemma}

\begin{proof}
Each factor $(v_i - v_j) = -(v_j - v_i)$, and there are $\binom{m}{2} = m(m-1)/2$
such factors.
\end{proof}

\begin{lemma}
\label{lem.det.vander.sign-Fin-revPerm}
\lean{AlgebraicCombinatorics.Determinants.sign_Fin_revPerm}
\leanhelper
\leanok
The sign of the column reversal permutation $\tau$ on $\{1, 2, \ldots, m\}$
(sending $j$ to $m + 1 - j$) satisfies
\[
  \operatorname{sign}(\tau) = (-1)^{m(m-1)/2}.
\]
\end{lemma}

\begin{proof}
For $i < j$, we have $\tau(i) > \tau(j)$,
so every pair is an inversion. The total number of inversions is $\binom{m}{2}$.
\end{proof}

\begin{lemma}
\label{lem.det.vander.col-perm}
\lean{AlgebraicCombinatorics.Determinants.det_submatrix_col_perm'}
\leanhelper
\leanok
Let $A$ be an $m \times m$ matrix and $\sigma$ a permutation of $[m]$.
Then $\det(A \circ \sigma) = \operatorname{sign}(\sigma) \cdot \det A$,
where $A \circ \sigma$ denotes the column permutation $A_{i,\sigma(j)}$.
\end{lemma}

\begin{proof}
\leanok
Transpose to a row permutation, apply the row permutation determinant formula,
then transpose back.
\end{proof}

\begin{lemma}
\label{lem.det.vander.det-vandermondePolyMat}
\lean{AlgebraicCombinatorics.Determinants.det_vandermondePolyMat}
\leanhelper
\leanok
The determinant of the polynomial Vandermonde matrix (with entries $x_i^{m-1-j}$)
equals
\[
  \det\!\left(\left(x_i^{m-1-j}\right)_{i,j}\right) =
  \operatorname{sign}(\tau) \cdot
  \det\!\left(\left(x_i^{j}\right)_{i,j}\right),
\]
where $\tau$ is the column reversal permutation (sending $j$ to $m - 1 - j$).
This relates the textbook matrix convention to Mathlib's via column reversal.
\end{lemma}

\begin{proof}
The polynomial Vandermonde matrix is Mathlib's Vandermonde matrix with columns
permuted by the reversal permutation, so the determinant picks up the sign factor.
\end{proof}

\subsection{Helpers for the Desnanot--Jacobi proof}

\begin{lemma}
\label{lem.det.adj-corners}
\lean{AlgebraicCombinatorics.Determinants.det_adjugate_corners}
\leanhelper
\leanok
Let $A \in K^{n \times n}$ with $n \geq 2$. Then
\begin{align*}
  & \det(A_{\sim 1, \sim 1}) \cdot \det(A_{\sim n, \sim n})
    - \det(A_{\sim 1, \sim n}) \cdot \det(A_{\sim n, \sim 1}) \\
  &= (\operatorname{adj} A)_{0,0} \cdot (\operatorname{adj} A)_{n-1,n-1}
     - (\operatorname{adj} A)_{0,n-1} \cdot (\operatorname{adj} A)_{n-1,0}.
\end{align*}
\end{lemma}

\begin{proof}
Direct computation using $(\operatorname{adj} A)_{i,j} = (-1)^{i+j} \det(A_{\sim j, \sim i})$.
\end{proof}

\begin{lemma}
\label{lem.det.desnanot-jacobi-base}
\lean{AlgebraicCombinatorics.Determinants.desnanot_jacobi_base}
\leanhelper
\leanok
The Desnanot--Jacobi identity holds for $2 \times 2$ matrices by direct
computation.
\end{lemma}

\begin{proof}
\leanok
The identity $\det A \cdot \det(A') = \det(A_{\sim 1,\sim 1}) \cdot \det(A_{\sim 2,\sim 2}) - \det(A_{\sim 1,\sim 2}) \cdot \det(A_{\sim 2,\sim 1})$
reduces to $\det A \cdot 1 = A_{2,2} \cdot A_{1,1} - A_{1,2} \cdot A_{2,1}$
for a $2 \times 2$ matrix, which is the definition of $\det A$.
\end{proof}

\begin{lemma}
\label{lem.det.jacobi-complement-2x2-base}
\lean{AlgebraicCombinatorics.Determinants.jacobi_complement_2x2_base}
\leanhelper
\leanok
For a $2 \times 2$ matrix $A$, the $2 \times 2$ determinant of the corner
entries of $\operatorname{adj} A$ satisfies
\[
  (\operatorname{adj} A)_{0,0} \cdot (\operatorname{adj} A)_{1,1}
  - (\operatorname{adj} A)_{0,1} \cdot (\operatorname{adj} A)_{1,0}
  = \det A.
\]
This is the base case of the Jacobi complementary minor theorem
(since $\det(A') = 1$ for the $0 \times 0$ inner submatrix).
\end{lemma}

\begin{proof}
\leanok
Direct computation using the explicit $2 \times 2$ adjugate formula.
\end{proof}

\begin{lemma}
\label{lem.det.jacobi-complement-2x2-three}
\lean{AlgebraicCombinatorics.Determinants.jacobi_complement_2x2_three}
\leanhelper
\leanok
For a $3 \times 3$ matrix $A$, the $2 \times 2$ determinant of the corner
entries of $\operatorname{adj} A$ at positions $\{0, 2\} \times \{0, 2\}$ satisfies
\[
  (\operatorname{adj} A)_{0,0} \cdot (\operatorname{adj} A)_{2,2}
  - (\operatorname{adj} A)_{0,2} \cdot (\operatorname{adj} A)_{2,0}
  = \det A \cdot A_{1,1}.
\]
This verifies the Jacobi complementary minor theorem for $3 \times 3$
matrices with $P = Q = \{0, 2\}$.
\end{lemma}

\begin{proof}
\leanok
Direct expansion of the adjugate entries and the $3 \times 3$ determinant,
followed by algebraic simplification.
\end{proof}

\subsection{Helpers for the generalized Desnanot--Jacobi proof}

\begin{convention}
\label{conv.det.submatrixRemove2}
\lean{AlgebraicCombinatorics.Determinants.submatrixRemove2}
\leanok
Let $n \in \mathbb{N}$ with $n \geq 2$. Let $A$ be an $n \times n$-matrix.
For $p < q$ in $[n]$ and $u < v$ in $[n]$, define
\[
  \operatorname{sub}_{[n] \setminus \{p,q\}}^{[n] \setminus \{u,v\}} A
\]
as the $(n-2) \times (n-2)$-submatrix of $A$ obtained by removing rows $p, q$
and columns $u, v$. The row and column indexing uses
the order-preserving injection $[n-2] \hookrightarrow [n]$
that skips the two removed indices in order.
\end{convention}

\begin{lemma}
\label{lem.det.submatrixRemove2-eq-submatrixDet}
\lean{AlgebraicCombinatorics.Determinants.submatrixRemove2_det_eq_submatrixDet_compl}
\leanhelper
\leanok
Let $A$ be an $(m+2) \times (m+2)$ matrix, and let $p < q$ and $u < v$ be
indices. Then
\[
  \det\!\left(\operatorname{sub}_{[n]\setminus\{p,q\}}^{[n]\setminus\{u,v\}} A\right)
  = \det\!\left(\operatorname{sub}_{\{p,q\}^c}^{\{u,v\}^c} A\right),
\]
where the right-hand side uses the complement-based submatrix construction.
This connects the explicit order-preserving injection used in the submatrix
construction to the general complement-based definition, enabling the
use of Jacobi's complementary minor theorem.
\end{lemma}

\begin{proof}
\leanok
The proof shows that the order-preserving injection that skips $p$ and $q$,
being strictly monotone and having range equal to $\{p,q\}^c$, coincides with
the canonical order-preserving enumeration of the complement set, so the two
submatrix constructions yield the same matrix.
\end{proof}

\subsection{Helpers for the adjugate and inverse relationship}

\begin{lemma}
\label{lem.det.inv-eq-adjugate-div}
\lean{AlgebraicCombinatorics.Determinants.inv_apply_eq_adjugate_div_det}
\leanhelper
\leanok
Let $K$ be a field and $A \in K^{n \times n}$ with $\det A \neq 0$.
Then for all $i, j$,
\[
  A^{-1}_{i,j} = \frac{(\operatorname{adj} A)_{i,j}}{\det A}.
\]
\end{lemma}

\begin{proof}
Follows from $A^{-1} = (\det A)^{-1} \cdot \operatorname{adj} A$.
\end{proof}

\begin{lemma}
\label{lem.det.adjugate-submatrix-eq-smul}
\lean{AlgebraicCombinatorics.Determinants.adjugate_submatrix_eq_smul}
\leanhelper
\leanok
Let $A$ be an $n \times n$ matrix with $\det A$ a unit. Then for any
index functions $f$ and $g$,
\[
  (\operatorname{adj} A)_{f(i), g(j)}
  = \det A \cdot (A^{-1})_{f(i), g(j)}.
\]
That is, the adjugate submatrix equals $\det A$ times the corresponding
inverse submatrix.
\end{lemma}

\begin{proof}
Entry-wise, $(\operatorname{adj} A)_{i,j} = \det A \cdot A^{-1}_{i,j}$
for invertible $A$.
\end{proof}

\begin{lemma}
\label{lem.det.det-adjugate-submatrix-finset}
\lean{AlgebraicCombinatorics.Determinants.det_adjugate_submatrix_finset}
\leanhelper
\leanok
Let $A$ be an $m \times m$ matrix with $\det A$ a unit.
Let $P, Q \subseteq [m]$ with $|P| = |Q| = k$. Then
\[
  \det\!\left(\operatorname{sub}_P^Q(\operatorname{adj} A)\right)
  = (\det A)^k \cdot \det\!\left(\operatorname{sub}_P^Q(A^{-1})\right).
\]
This reduces Jacobi's complementary minor theorem to the complementary
minor theorem for inverse matrices.
\end{lemma}

\begin{proof}
By Lemma~\ref{lem.det.adjugate-submatrix-eq-smul}, the adjugate submatrix
is $\det A$ times the inverse submatrix. The determinant of a scalar
multiple is $(\det A)^k$ times the original determinant.
\end{proof}

\subsection{Helpers for block matrix decomposition}

\begin{lemma}
\label{lem.det.complementary-minor-block}
\lean{AlgebraicCombinatorics.Determinants.complementary_minor_block}
\leanhelper
\leanok
For a block matrix $M = \begin{pmatrix} A & B \\ C & D \end{pmatrix}$ with $D$
invertible and the Schur complement $A - BD^{-1}C$ invertible, we have
\[
  \det(M) \cdot \det\!\left((M^{-1})_{\text{top-left}}\right) = \det D.
\]
\end{lemma}

\begin{proof}
By the Schur complement formula, $\det(M) = \det(D) \cdot \det(A - BD^{-1}C)$.
The top-left block of $M^{-1}$ is $(A - BD^{-1}C)^{-1}$, whose determinant
is $\det(A - BD^{-1}C)^{-1}$. Multiplying gives $\det(D)$.
\end{proof}

\begin{lemma}
\label{lem.det.adjugate-toBlocks11}
\lean{AlgebraicCombinatorics.Determinants.adjugate_toBlocks₁₁_det_of_invertible}
\leanhelper
\leanok
Let $M = \begin{pmatrix} A & B \\ C & D \end{pmatrix}$ be an invertible
block matrix with $D$ invertible and Schur complement $A - BD^{-1}C$
invertible. Let $k = |A|$ (the size of the top-left block) with $k \geq 1$.
Then
\[
  \det\!\left((\operatorname{adj} M)_{\text{top-left}}\right)
  = (\det M)^{k-1} \cdot \det D.
\]
\end{lemma}

\begin{proof}
Since $\operatorname{adj}(M) = \det(M) \cdot M^{-1}$ for invertible $M$,
the top-left block of $\operatorname{adj}(M)$ is
$\det(M) \cdot (M^{-1})_{\text{top-left}}$. Taking the $k \times k$
determinant yields $(\det M)^k \cdot \det((M^{-1})_{\text{top-left}})$.
By Lemma~\ref{lem.det.complementary-minor-block},
$\det(M) \cdot \det((M^{-1})_{\text{top-left}}) = \det D$,
so $\det((M^{-1})_{\text{top-left}}) = \det D / \det M$, giving the result.
\end{proof}

\begin{definition}
\label{def.det.sortEquivPQ}
\lean{AlgebraicCombinatorics.Determinants.sortEquivPQ}
\leanhelper
\leanok
Given a subset $P$ of $[m]$, the \emph{sorting equivalence} is a bijection
\[
  \{1, \ldots, |P|\} \sqcup \{1, \ldots, m - |P|\} \xrightarrow{\sim} [m]
\]
that sends the $i$-th element of the left component to the $i$-th smallest element of $P$
and the $j$-th element of the right component to the $j$-th smallest element of $P^c$.
This equivalence is used to decompose matrices into block form with respect to the partition
$[m] = P \sqcup P^c$.
\end{definition}

\begin{lemma}
\label{lem.det.inv-submatrix-equiv}
\lean{AlgebraicCombinatorics.Determinants.inv_submatrix_equiv}
\leanhelper
\leanok
Let $K$ be a field, $A \in K^{m \times m}$ with $\det A \neq 0$, and
$\sigma, \tau$ bijections from $[n]$ to $[m]$. Then
\[
  (A_{\sigma, \tau})^{-1} = (A^{-1})_{\tau, \sigma},
\]
where subscripts denote the submatrix obtained by reindexing rows and columns
via the given bijections.
\end{lemma}

\begin{proof}
Verify directly that $A_{\sigma,\tau} \cdot (A^{-1})_{\tau,\sigma} = I$ using
$(A \cdot A^{-1})_{\sigma(i), \sigma(j)} = \delta_{ij}$.
\end{proof}

\begin{lemma}
\label{lem.det.det-submatrix-equiv}
\lean{AlgebraicCombinatorics.Determinants.det_submatrix_equiv_equiv}
\leanhelper
\leanok
Let $A \in K^{m \times m}$ and $f, g : [n] \xrightarrow{\sim} [m]$
be bijections. Then
\[
  \det(A_{f, g}) = \operatorname{sign}(g^{-1} \circ f) \cdot \det A.
\]
\end{lemma}

\begin{proof}
Rewrite the submatrix as a reindexed matrix and apply the standard
determinant reindexing formula.
\end{proof}

\subsection{Helpers for the complementary minor theorem proof}

\begin{lemma}
\label{lem.det.complementary-minor-inverse}
\lean{AlgebraicCombinatorics.Determinants.complementary_minor_inverse}
\leanhelper
\leanok
Let $K$ be a field. Let $A \in K^{m \times m}$ be an invertible matrix.
Let $P, Q \subseteq [m]$ with $|P| = |Q|$. Then
\[
  \det A \cdot \det(\operatorname{sub}_P^Q(A^{-1}))
  = (-1)^{\operatorname{sum} P + \operatorname{sum} Q}
    \det(\operatorname{sub}_{\widetilde{Q}}^{\widetilde{P}} A).
\]
\end{lemma}

\begin{proof}
The proof uses block matrix decomposition. Define sorting equivalences
$\sigma$ and $\tau$ for $Q$ and $P$ respectively
(Definition~\ref{def.det.sortEquivPQ}), which reorder indices so that
$P$-elements come first. The permuted matrix $M = A_{\sigma, \tau}$
decomposes into blocks: $M_{\text{top-left}} = \operatorname{sub}_Q^P(A)$
and $M_{\text{bottom-right}} = \operatorname{sub}_{Q^c}^{P^c}(A)$.
By Lemma~\ref{lem.det.inv-submatrix-equiv},
$(M^{-1})_{\text{top-left}} = \operatorname{sub}_P^Q(A^{-1})$.
By Lemma~\ref{lem.det.complementary-minor-block},
$\det(M) \cdot \det((M^{-1})_{\text{top-left}}) = \det(M_{\text{bottom-right}})$.
The sign factor $(-1)^{\operatorname{sum} P + \operatorname{sum} Q}$
arises from the determinant of the sorting permutation
(Lemma~\ref{lem.det.det-submatrix-equiv}).
\end{proof}

\begin{lemma}
\label{lem.det.jacobi-field}
\lean{AlgebraicCombinatorics.Determinants.jacobi_complementary_minor_field}
\leanhelper
\leanok
Let $K$ be a field. Let $A \in K^{m \times m}$ with $\det A \neq 0$.
Let $P, Q \subseteq [m]$ with $|P| = |Q| \geq 1$. Then
\[
  \det(\operatorname{sub}_P^Q(\operatorname{adj} A))
  = (-1)^{\operatorname{sum} P + \operatorname{sum} Q}
    (\det A)^{|Q|-1}
    \det(\operatorname{sub}_{\widetilde{Q}}^{\widetilde{P}} A).
\]
\end{lemma}

\begin{proof}
By Lemma~\ref{lem.det.det-adjugate-submatrix-finset},
\[
  \det(\operatorname{sub}_P^Q(\operatorname{adj} A))
  = (\det A)^{|P|} \cdot \det(\operatorname{sub}_P^Q(A^{-1})).
\]
By Lemma~\ref{lem.det.complementary-minor-inverse},
\[
  \det A \cdot \det(\operatorname{sub}_P^Q(A^{-1}))
  = (-1)^{\operatorname{sum} P + \operatorname{sum} Q}
    \det(\operatorname{sub}_{\widetilde{Q}}^{\widetilde{P}} A),
\]
so $\det(\operatorname{sub}_P^Q(A^{-1})) = (-1)^{\operatorname{sum} P + \operatorname{sum} Q}
\det(\operatorname{sub}_{\widetilde{Q}}^{\widetilde{P}} A) / \det A$.
Since $|P| = |Q|$, the exponent becomes $|Q| - 1$.
\end{proof}

\subsection{Helpers for the Cauchy determinant proof (factor hunting)}

\begin{definition}
\label{def.det.substXiToXj}
\lean{AlgebraicCombinatorics.Determinants.substXiToXj}
\leanhelper
\leanok
The \emph{variable substitution} $\varphi_{i \to j}$ is the
algebra homomorphism
$R[x_1, \ldots, x_n] \to R[x_1, \ldots, x_n]$ that sends $x_i \mapsto x_j$
and fixes all other variables.
\end{definition}

\begin{lemma}
\label{lem.det.X-sub-X-dvd-sub-subst}
\lean{AlgebraicCombinatorics.Determinants.X_sub_X_dvd_sub_subst}
\leanhelper
\leanok
Let $p \in R[x_1, \ldots, x_n]$ be a multivariate polynomial and $i \neq j$.
Then $(x_i - x_j) \mid p - p|_{x_i = x_j}$, where $p|_{x_i = x_j}$
denotes the substitution of $x_j$ for $x_i$.
\end{lemma}

\begin{proof}
\leanok
By induction on the structure of $p$ (using the polynomial recursor on
constants, sums, and products by a variable).
The base case (constants) is trivial. The additive case follows from
divisibility of sums. The multiplicative case $p \cdot x_k$ splits:
if $k = i$, factor as $(p - p|_{x_i = x_j}) \cdot x_i + p|_{x_i = x_j} \cdot (x_i - x_j)$;
otherwise, factor as $(p - p|_{x_i = x_j}) \cdot x_k$.
\end{proof}

\begin{lemma}
\label{lem.det.X-sub-X-dvd-eval-zero}
\lean{AlgebraicCombinatorics.Determinants.X_sub_X_dvd_of_eval₂_eq_zero_fin_01}
\leanhelper
\leanok
Let $R$ be an integral domain, and let
$P \in R[x_0, x_1, \ldots, x_{m+1}]$. If $P$ vanishes when $x_0$ is
replaced by $x_1$ (i.e., $P|_{x_0 = x_1} = 0$), then
$(x_0 - x_1) \mid P$.
\end{lemma}

\begin{proof}
\leanok
View $P$ as a univariate polynomial in $x_0$
over $R[x_1, \ldots, x_{m+1}]$ (using the canonical isomorphism
between $R[x_0, x_1, \ldots, x_{m+1}]$ and $R[x_1, \ldots, x_{m+1}][x_0]$).
The hypothesis says that $x_1$ is a root
of this univariate polynomial, so $(X - x_1)$ divides it by the polynomial
root theorem. Pulling back through the isomorphism yields
$(x_0 - x_1) \mid P$.
\end{proof}

\begin{lemma}
\label{lem.det.X-sub-X-irreducible}
\lean{AlgebraicCombinatorics.Determinants.X_sub_X_irreducible}
\leanhelper
\leanok
For distinct variables $x_i, x_j$ in $\mathbb{Z}[x_1, \ldots, x_n]$,
the polynomial $x_i - x_j$ is irreducible.
\end{lemma}

\begin{proof}
\leanok
Since $x_i - x_j$ has total degree~$1$ and is primitive (its content is~$1$),
it is irreducible. The degree-$1$ argument: any nontrivial factorization
would require one factor of degree~$0$ (a constant), but a primitive
polynomial of positive degree cannot have a non-unit constant factor.
\end{proof}

\begin{lemma}
\label{lem.det.X-sub-X-totalDegree}
\lean{AlgebraicCombinatorics.Determinants.X_sub_X_totalDegree_eq_one}
\leanhelper
\leanok
For distinct $i, j$, the total degree of $x_i - x_j$ in
$\mathbb{Z}[x_1, \ldots, x_n]$ is~$1$.
\end{lemma}

\begin{proof}
\leanok
The upper bound follows from $\deg(f - g) \leq \max(\deg f, \deg g)$
and $\deg(X_i) = 1$. The lower bound follows from the monomial
$X_i^1$ having nonzero coefficient in $x_i - x_j$.
\end{proof}

\begin{lemma}
\label{lem.det.X-sub-X-isPrimitive}
\lean{AlgebraicCombinatorics.Determinants.X_sub_X_isPrimitive}
\leanhelper
\leanok
For distinct $i, j$, the polynomial $x_i - x_j \in \mathbb{Z}[x_1, \ldots, x_n]$
is primitive (its content is~$1$).
\end{lemma}

\begin{proof}
\leanok
The coefficient of the monomial $x_i^1$ is $1$, so the GCD of the
coefficients is $1$.
\end{proof}

\begin{lemma}
\label{lem.det.X-sub-X-isRelPrime}
\lean{AlgebraicCombinatorics.Determinants.X_sub_X_isRelPrime}
\leanhelper
\leanok
Let $i \neq j$ and $k \neq l$ be pairs of distinct indices such that
$\{i,j\}$ and $\{k,l\}$ are not equal as unordered pairs. Then
$(x_i - x_j)$ and $(x_k - x_l)$ are coprime in
$\mathbb{Z}[x_1, \ldots, x_n]$.
\end{lemma}

\begin{proof}
\leanok
Both are irreducible by Lemma~\ref{lem.det.X-sub-X-irreducible}, and they
are non-associate (not equal up to unit multiple) since they involve
different variable pairs. Distinct irreducibles in a UFD are coprime.
\end{proof}

\begin{lemma}
\label{lem.det.cauchy-poly-dvd-x}
\lean{AlgebraicCombinatorics.Determinants.polyCauchyMat'_det_dvd_x_sub_x}
\leanhelper
\leanok
Let $C(m)$ be the $m \times m$ matrix with entry
$(i,j) = \prod_{k \neq j}(X_i + Y_k)$ over
$\mathbb{Z}[X_1, \ldots, X_m, Y_1, \ldots, Y_m]$.
For distinct $i, j$, we have $(X_i - X_j) \mid \det(C(m))$.
\end{lemma}

\begin{proof}
\leanok
Substituting $X_j$ for $X_i$ makes rows $i$ and $j$ identical, so the
determinant vanishes. By Lemma~\ref{lem.det.X-sub-X-dvd-eval-zero},
$(X_i - X_j)$ divides the determinant.
\end{proof}

\begin{lemma}
\label{lem.det.cauchy-poly-dvd-y}
\lean{AlgebraicCombinatorics.Determinants.polyCauchyMat'_det_dvd_y_sub_y}
\leanhelper
\leanok
For distinct $i, j$, we have
$(Y_i - Y_j) \mid \det(C(m))$,
where $C(m)$ is the polynomial Cauchy matrix from Lemma~\ref{lem.det.cauchy-poly-dvd-x}.
\end{lemma}

\begin{proof}
\leanok
Substituting $Y_j$ for $Y_i$ makes columns $i$ and $j$ proportional
(both become $\prod_{k \neq i, k \neq j}(X_r + Y_k) \cdot (X_r + Y_j)$),
so the determinant vanishes. The divisibility follows by
Lemma~\ref{lem.det.X-sub-X-dvd-eval-zero}.
\end{proof}

\begin{lemma}
\label{lem.det.cauchy-det-of-poly}
\lean{AlgebraicCombinatorics.Determinants.cauchy_det_of_poly}
\leanhelper
\leanok
Let $K$ be a field. Let $x_1, \ldots, x_m, y_1, \ldots, y_m \in K$ with
$x_i + y_j \neq 0$ for all $i, j$. Then
\[
  \det\!\left(\frac{1}{x_i + y_j}\right)_{i,j}
  = \frac{\prod_{i < j}(x_i - x_j)(y_i - y_j)}{\prod_{i,j}(x_i + y_j)}.
\]
\end{lemma}

\begin{proof}
The polynomial Cauchy identity (Lemma~\ref{lem.det.cauchy-poly}) gives
$\det(\prod_{k \neq j}(x_i + y_k)) = \prod_{i < j}(x_i - x_j)(y_i - y_j)$.
The left-hand side equals $\prod_{i,j}(x_i + y_j) \cdot \det(1/(x_i + y_j))$
since each row of the Cauchy matrix can be multiplied by
$\prod_j (x_i + y_j)$ to clear denominators. Dividing gives the result.
\end{proof}
